% =============================================================================
%  esprit-style.tex — Habillage graphique du gabarit ESPRIT
%  Chargé par main.tex après les paquets.
% =============================================================================

% ── Palette ESPRIT ───────────────────────────────────────────────────────────
\definecolor{espritred}{HTML}{CE1126}
\definecolor{espritdark}{HTML}{3C3C3B}
\definecolor{espritgrey}{HTML}{9C9B9B}
\definecolor{espritlight}{HTML}{D9D9D9}
\definecolor{esprittable}{HTML}{1F3864}   % bandeau d'en-tete des tableaux
\definecolor{espritband}{HTML}{C00000}

% ── Motif de triangles : coin superieur droit ────────────────────────────────
\newcommand{\espritTrianglesHaut}{%
  \begin{tikzpicture}[x=1mm,y=1mm]
    % rangee de petits triangles, du plus clair au plus sature
    \fill[espritlight]  (0,0)   -- (5,0)   -- (2.5,4.5)  -- cycle;
    \fill[espritgrey]   (6,1)   -- (12,1)  -- (9,6)      -- cycle;
    \fill[espritred]    (13,0)  -- (20,0)  -- (16.5,6)   -- cycle;
    \fill[espritlight]  (21,2)  -- (26,2)  -- (23.5,6.5) -- cycle;
    \fill[espritdark]   (27,0)  -- (34,0)  -- (30.5,6)   -- cycle;
    \fill[espritgrey]   (35,1)  -- (40,1)  -- (37.5,5.5) -- cycle;
    \fill[espritred]    (41,0)  -- (49,0)  -- (45,7)     -- cycle;
    \fill[espritdark]   (50,2)  -- (56,2)  -- (53,7)     -- cycle;
    \fill[espritgrey]   (57,0)  -- (64,0)  -- (60.5,6)   -- cycle;
    % seconde rangee, decalee
    \fill[espritred]    (10,7)  -- (15,7)  -- (12.5,11)  -- cycle;
    \fill[espritlight]  (30,8)  -- (35,8)  -- (32.5,12)  -- cycle;
    \fill[espritdark]   (44,8)  -- (50,8)  -- (47,13)    -- cycle;
    \fill[espritgrey]   (56,7)  -- (62,7)  -- (59,12)    -- cycle;
  \end{tikzpicture}}

% ── Motif de triangles : bas de page ─────────────────────────────────────────
\newcommand{\espritTrianglesBas}{%
  \begin{tikzpicture}[x=1mm,y=1mm]
    \fill[espritred] (0,0)  -- (18,0) -- (9,13)  -- cycle;
    \fill[espritred] (34,0) -- (50,0) -- (42,12) -- cycle;
    \fill[espritred] (66,0) -- (82,0) -- (74,12) -- cycle;
  \end{tikzpicture}}

% ── Fond de page applique a toutes les pages du corps ────────────────────────
\newcommand{\espritHabillage}{%
  \AddToShipoutPictureBG{%
    % logo en haut a gauche
    \AtPageUpperLeft{\raisebox{-16mm}{\hspace*{14mm}%
      \includegraphics[height=9mm]{images/logo_esprit.png}}}%
    % triangles en haut a droite
    \AtPageUpperLeft{\raisebox{-20mm}{\hspace*{110mm}\espritTrianglesHaut}}%
    % triangles en bas
    \AtPageLowerLeft{\raisebox{4mm}{\hspace*{18mm}\espritTrianglesBas}}%
  }}

% ── Page de garde ────────────────────────────────────────────────────────────
\newcommand{\espritPageDeGarde}[7]{%
  % #1 annee  #2 specialite  #3 nom du produit  #4 sous-titre
  % #5 etudiant  #6 encadrant esprit  #7 encadrants entreprise
  \begin{titlepage}
  \thispagestyle{empty}
  \AddToShipoutPictureFG*{%
    % bande laterale gauche : triangles decoratifs
    \AtPageLowerLeft{\begin{tikzpicture}[x=1mm,y=1mm,remember picture,overlay]
      \fill[espritlight] (6,250)  -- (30,250) -- (18,228) -- cycle;
      \fill[espritred]   (10,262) -- (26,262) -- (18,248) -- cycle;
      \fill[espritgrey]  (2,215)  -- (22,215) -- (12,197) -- cycle;
      \fill[espritdark]  (14,190) -- (34,190) -- (24,172) -- cycle;
      \fill[espritlight] (4,160)  -- (26,160) -- (15,140) -- cycle;
      \fill[espritred]   (8,120)  -- (30,120) -- (19,100) -- cycle;
      \fill[espritgrey]  (2,84)   -- (20,84)  -- (11,68)  -- cycle;
      \fill[espritdark]  (16,60)  -- (36,60)  -- (26,42)  -- cycle;
      \fill[espritlight] (6,34)   -- (24,34)  -- (15,18)  -- cycle;
      % triangles disperses a droite
      \fill[espritlight] (150,240) -- (168,240) -- (159,224) -- cycle;
      \fill[espritgrey]  (172,214) -- (188,214) -- (180,200) -- cycle;
      \fill[espritlight] (140,120) -- (162,120) -- (151,100) -- cycle;
      \fill[espritgrey]  (176,86)  -- (192,86)  -- (184,72)  -- cycle;
    \end{tikzpicture}}}

  \centering
  \vspace*{4mm}
  \includegraphics[width=0.42\textwidth]{images/logo_esprit.png}\\[6mm]

  {\LARGE\bfseries\color{espritred} #1}\\[3mm]
  {\Huge\bfseries\color{espritdark} PROJET DE FIN D'ÉTUDES}\\[6mm]

  \colorbox{espritband}{\parbox{0.96\textwidth}{\centering\vspace{2mm}
    {\LARGE\bfseries\color{white} DIPLÔME NATIONAL D'INGÉNIEUR}\vspace{2mm}}}\\[8mm]

  {\large\bfseries\color{espritdark} SPÉCIALITÉ : #2}\\[10mm]

  {\Huge\bfseries #3}\\[5mm]
  {\large\itshape #4}\\[16mm]

  \begin{minipage}{0.82\textwidth}
    \raggedright\large
    \textbf{Réalisé par :} #5\\[2mm]
    \textbf{Encadrant ESPRIT :} #6\\[2mm]
    \textbf{Encadrants entreprise :} #7
  \end{minipage}\\[10mm]

  \includegraphics[width=0.16\textwidth]{images/cni_logo.jpg}

  \vfill
  {\color{espritgrey}\rule{0.9\textwidth}{0.6pt}}\\[3mm]
  {\footnotesize\color{espritgrey} Centre National de l'Informatique --- Tunis}
  \end{titlepage}
  \clearpage}

% ── Style des titres ─────────────────────────────────────────────────────────
\titleformat{\chapter}[hang]
  {\Large\bfseries}
  {\thechapter}{1em}{}
  [\vspace{-4pt}{\color{espritdark}\rule{\textwidth}{0.8pt}}]
\titlespacing*{\chapter}{0pt}{6pt}{18pt}

\titleformat{name=\chapter,numberless}[hang]
  {\Large\bfseries}{}{0pt}{}
  [\vspace{-4pt}{\color{espritdark}\rule{\textwidth}{0.8pt}}]

\titleformat{\section}{\large\bfseries}{\thesection}{0.8em}{}
\titlespacing*{\section}{0pt}{14pt}{6pt}

\titleformat{\subsection}{\normalsize\bfseries}{\thesubsection}{0.7em}{}
\titlespacing*{\subsection}{0pt}{11pt}{4pt}

\titleformat{\subsubsection}{\small\bfseries}{\thesubsubsection}{0.6em}{}
\titlespacing*{\subsubsection}{0pt}{9pt}{3pt}

% Titre de section non numerote employe dans les chapitres de sprint
\newcommand{\sectionsansnum}[1]{%
  \par\vspace{12pt}\noindent{\large\bfseries #1}\par\vspace{5pt}}

% ── Numerotation continue des figures et tableaux ────────────────────────────
\counterwithout{figure}{chapter}
\counterwithout{table}{chapter}

% ── Legendes ─────────────────────────────────────────────────────────────────
\captionsetup[figure]{
  labelfont     = {it,small},
  textfont      = {it,small},
  labelsep      = space,
  justification = centering,
  skip          = 4pt,
}
\captionsetup[table]{
  labelfont     = {small},
  textfont      = {small},
  labelsep      = space,
  justification = centering,
  position      = below,
  skip          = 4pt,
}
\captionsetup[subfigure]{labelfont={it,footnotesize},textfont={it,footnotesize}}

% ── Tableaux : bandeau d'en-tete bleu fonce, texte blanc ─────────────────────
\newcommand{\entete}[1]{\textcolor{white}{\textbf{#1}}}
\newcommand{\ligneentete}{\rowcolor{esprittable}}

% Tableau de description textuelle d'un cas d'utilisation
\newcommand{\legendecasu}{}
\newenvironment{casutilisation}[1]
  {\renewcommand{\legendecasu}{#1}%
   \begin{table}[H]\centering\small
   \renewcommand{\arraystretch}{1.25}
   \begin{tabular}{@{}>{\columncolor{esprittable}\color{white}\bfseries}p{3.3cm}
                     |p{10.4cm}@{}}
   \hline}
  {\\\hline\end{tabular}\par\vspace{2pt}\caption{\legendecasu}\end{table}}

% ── Boite « en cours de redaction » ──────────────────────────────────────────
\newcommand{\enattente}[1]{%
  \begin{center}
  \setlength{\fboxrule}{0.7pt}\setlength{\fboxsep}{8pt}%
  \fcolorbox{espritgrey}{white}{\parbox{0.84\textwidth}{\centering
    \color{espritdark}\small\itshape Section en cours de rédaction\\[5pt]
    \upshape\footnotesize #1}}%
  \end{center}}

% Ligne d'intitule d'epic dans le tableau du backlog
\newcommand{\epictitre}[1]{%
  \midrule
  \multicolumn{5}{@{}l}{\cellcolor{hypnolight}\textbf{#1}}\\
  \midrule}
